\section{Method}

\subsection{Overview and Problem Formulation}
Given a video $\mathcal{V}=\{I_t\}_{t=1}^{T}$ and a referring expression $q$, referring video object segmentation (RVOS) predicts a mask sequence $\mathcal{Y}=\{Y_t\}_{t=1}^{T}$ for the referred object(s). CRF-SAM3 factorizes this mapping into three explicit stages:
\begin{equation}
(q,\mathcal{V})
\xrightarrow{\text{Ground}}
\hat c
\xrightarrow[\mathcal{V}]{\text{Recall}}
\mathcal{C}(q)
\xrightarrow[\,q\,]{\text{Verify}}
\hat{\mathcal{S}}\subseteq\mathcal{C}(q),
\end{equation}
where $\hat c$ is a compact, referent-aware recall concept, $\mathcal{C}(q)=\{\tau_i\}_{i=1}^{N_c}$ is a set of candidate object tubes, and $\hat{\mathcal{S}}$ is the set of selected tubes whose union forms the prediction. The three stages solve deliberately different subproblems. Ground asks which visual concept should be used to search the video; Recall asks which temporally coherent object hypotheses satisfy that concept; Verify asks which recalled hypotheses satisfy the complete referring expression.

This factorization is intentionally asymmetric. The first two stages are recall-oriented: it is preferable to preserve several plausible same-category tubes than to discard the true referent too early. Fine-grained attributes, actions, relations, and temporal clauses are therefore not forced into the SAM3 prompt. Instead, they are retained in the original expression and used again during candidate verification. The Ground module and the verifier are trained separately (Secs.~\ref{sec:ground} and \ref{sec:verify}); SAM3 is frozen and no gradient is propagated through candidate generation. Figure~\ref{fig:framework} summarizes this data flow.

\subsection{Referent-Aware Concept Grounding}
\label{sec:ground}
\paragraph{Visual and language representations.}
We sample $M$ RGB frames from the video and encode them with the frozen Qwen3-VL vision encoder. Let $V=\{v_n\}_{n=1}^{N_v}$ denote the resulting merged visual tokens, i.e., the embeddings that Qwen3-VL inserts at its image placeholder positions. We do not crop the input around the target, nor do we remove visual tokens. In parallel, the expression $q$ is tokenized on its own and embedded with the frozen Qwen3-VL token-embedding table. A lightweight expression conditioner projects these embeddings to the query dimension, prepends a learnable summary token, and applies a small bidirectional Transformer encoder, producing conditioned text features $T$.

The two streams play different roles: $V$ preserves the original RGB context, whereas $T$ summarizes the complete referring expression for target localization. The Qwen3-VL backbone itself, including its vision encoder and language decoder, remains frozen throughout this stage.

\paragraph{Category-agnostic target-query grounding.}
We introduce $K$ learnable query embeddings $Q^0=\{q_k^0\}_{k=1}^{K}$. They are not tied to semantic categories; each slot is free to represent a possible expression-conditioned target hypothesis. Let $\bar V=P_v(V)$ denote the visual tokens linearly projected to the query dimension. A query decoder of $L$ pre-norm residual layers updates the slots by query self-attention, cross-attention to the text features, cross-attention to the visual tokens, and an FFN:
\begin{equation}
Q^{\ell}=\mathcal{D}_{\ell}\!\left(Q^{\ell-1},T,\bar V\right),
\qquad \ell=1,\ldots,L .
\end{equation}
The slots therefore use the entire expression while grounding themselves directly in the original visual-token sequence.

The final query states are converted to token responses with a scaled cosine similarity:
\begin{equation}
R_{k,n}
=
\gamma\,
\left\langle
\operatorname{norm}(W_q q_k^{L}),
\operatorname{norm}(W_v \bar v_n)
\right\rangle,
\end{equation}
where $\gamma$ is a fixed response scale. The response matrix $R$ is an explicit intermediate representation of where each slot finds evidence for the referred target.

\paragraph{Mask supervision without mask input.}
Ground-truth masks supervise these token responses during training only. They are never rendered into the RGB inputs and are not supplied to the model at inference. For each training frame, the binary target mask is resized to the patch grid of the vision encoder and average-pooled over each spatial-merge window; a merged token is labelled positive if the target covers at least a fraction $\rho$ of its window, giving token labels $y\in\{0,1\}^{N_v}$. We compute a matching cost for every slot,
\begin{equation}
\ell_k^{\mathrm{mask}}
=
\mathcal{L}_{\mathrm{BCE}}(R_k,y)
+
\mathcal{L}_{\mathrm{Dice}}\!\left(\sigma(R_k),y\right),
\end{equation}
and assign the spatial supervision to the best-matching slot,
\begin{equation}
 k^{*}=\arg\min_k \ell_k^{\mathrm{mask}},
 \qquad
 \mathcal{L}_{\mathrm{ground}}
 =\ell_{k^{*}}^{\mathrm{mask}}.
\end{equation}
The best-slot rule is used only because the slots are exchangeable; it does not turn them into fixed semantic prototypes.

\paragraph{Soft visual modulation.}
The token responses are not used to crop the frame or prune the visual sequence. Instead, they provide a soft residual gate. During training, the selected slot produces
\begin{equation}
 m_n^{\mathrm{train}}=\sigma(R_{k^{*},n}),
\end{equation}
and a bottleneck adapter $A(\cdot)$ (LayerNorm, linear down-projection, SiLU, linear up-projection) modulates each original visual token as
\begin{equation}
 v'_n
 =
 v_n+\alpha\,m_n A(v_n),
\label{eq:soft_modulation}
\end{equation}
where $\alpha$ is a learnable scalar. The last layer of $A$ is zero-initialized, so the module starts as the identity and the frozen decoder initially receives unmodified visual tokens. The same gate and adapter are applied to Qwen3-VL's DeepStack visual embeddings, which are injected into intermediate decoder layers. Grounding thus changes the contribution of target-relevant evidence while preserving token count, order, and the surrounding scene context expected by the frozen decoder.

\paragraph{Concept generation.}
The modulated visual embeddings are written back to their original placeholder positions. The frozen Qwen3-VL decoder receives an instruction prompt containing the referring expression and the $M$ frames, and generates a short JSON answer \texttt{\{"category": $\hat c$\}}. During training, the target is a category label $c^{\star}$ for the referred object, and $\mathcal{L}_{\mathrm{cat}}$ is the token-level cross-entropy of this answer. The category labels are used only as training targets; their construction is described in Sec.~\ref{sec:setup}. The grounding objective is
\begin{equation}
\mathcal{L}_{G}
=
\lambda_g\mathcal{L}_{\mathrm{ground}}
+
\lambda_c\mathcal{L}_{\mathrm{cat}}.
\end{equation}
Because the gate $m^{\mathrm{train}}$ is not detached, $\mathcal{L}_{\mathrm{cat}}$ is back-propagated through the frozen decoder into the adapter and, through $m$, into the query decoder. The target queries are therefore shaped both by \emph{where} the referent is ($\mathcal{L}_{\mathrm{ground}}$) and by \emph{which concept} the decoder should emit from the modulated evidence ($\mathcal{L}_{\mathrm{cat}}$). Only the expression conditioner, query decoder, visual projection $P_v$, response head ($W_q,W_v$), and modulation adapter are optimized.

\paragraph{Mask-free inference.}
At test time, the ground-truth-based assignment $k^{*}$ is unavailable. Rather than committing to one slot, we aggregate the responses over all slots at every visual-token position:
\begin{equation}
 m_n^{\mathrm{test}}
 =
 \max_k \sigma(R_{k,n}).
\label{eq:test_query_aggregation}
\end{equation}
The resulting response is inserted into Eq.~\eqref{eq:soft_modulation}, after which the decoder generates $\hat c$ by greedy decoding. Inference therefore requires only the RGB frames and the referring expression. The training masks serve as privileged spatial supervision for learning \emph{where the expression refers}, not as an inference-time visual cue.

\subsection{SAM3 Candidate Recall}
\label{sec:recall}
The grounded concept $\hat c$ is an interface to object discovery rather than the final semantic answer. We provide $\hat c$ as the concept prompt, together with the RGB video, to SAM3~\cite{carion2025sam3}, using the training-free decoupled memory selection of SAM3-DMS~\cite{shen2026sam3dms} for multi-object propagation. Each retained object identity forms a candidate tube
\begin{equation}
\tau_i=\{M_t^{(i)}\}_{t=1}^{T},
\qquad
\mathcal{C}(q)=\{\tau_i\}_{i=1}^{N_c}.
\end{equation}
SAM3 first detects up to $K_r$ concept instances, ranked by its detection score, and only the top $K_{\mathrm{keep}}\le K_r$ identities are propagated and passed on, so $N_c\le K_{\mathrm{keep}}$. SAM3 is kept frozen. This stage is intentionally recall-oriented: multiple same-category instances may remain in $\mathcal{C}(q)$ even when the compact concept cannot distinguish them. Such ambiguity is acceptable as long as the true referent remains available to the verifier. If SAM3 returns no instance, $\mathcal{C}(q)=\emptyset$ and the prediction is empty.

\paragraph{Offline candidates and IoU labels.}
For verifier training, candidate tubes are generated offline with the same recall procedure and cached. Since the verifier only consumes candidate tubes, the training cache is built from training-time concepts rather than from the grounder output (Sec.~\ref{sec:setup}). Candidate supervision is obtained from overlap with the ground-truth tube rather than from the SAM3 score. We use the sequence-level region similarity
\begin{equation}
J_i
=
\frac{1}{|\Omega_G|}
\sum_{t\in\Omega_G}
\operatorname{IoU}\!\left(M_t^{(i)},G_t\right),
\end{equation}
where $\Omega_G$ denotes the annotated ground-truth frames and a candidate missing on such a frame contributes an empty mask. A candidate is positive if $J_i\ge\theta_{\mathrm{pos}}$ and negative if $J_i\le\theta_{\mathrm{neg}}$. These IoU labels supervise ranking only and are never required at inference.

\subsection{Candidate-Aware Verification}
\label{sec:verify}
\paragraph{Candidate-conditioned visual evidence.}
The compact recall concept is discarded once the candidate set has been constructed. Verification instead returns to the original expression $q$. For each candidate tube $\tau_i$, we uniformly sample $F$ frames from the frames on which the candidate is present and mark the candidate in each frame with a semi-transparent mask overlay. These frames form the visual evidence for that specific hypothesis. The verifier prompt instructs Qwen3-VL to treat the marked object as the anchor subject of the expression and to answer \emph{Yes} if the expression becomes true when grounded on that object in the frame sequence, and \emph{No} otherwise. This anchor formulation is important: without it, a multimodal model can answer whether the expression is true somewhere in the scene rather than whether it is true for the candidate under evaluation. At inference, when the expression contains a color word, the candidate is marked by its bounding box instead of a filled overlay so that the overlay does not alter the color evidence.

Let $\ell_i^{\mathrm{yes}}$ and $\ell_i^{\mathrm{no}}$ be the next-token logits of \emph{Yes} and \emph{No} after the prompt. The candidate score is
\begin{equation}
 s_i
 =
 P(\mathrm{Yes}\mid q,\tau_i)
 =\frac{\exp\ell_i^{\mathrm{yes}}}{\exp\ell_i^{\mathrm{yes}}+\exp\ell_i^{\mathrm{no}}}.
\end{equation}
The vision encoder of the verifier is frozen, and LoRA adapters on the attention and MLP projections of the language model adapt Qwen3-VL to candidate-level discrimination. The same score is used for training and for ranking at inference.

\paragraph{Listwise candidate competition.}
Scoring candidates independently does not reflect the decision made at inference, where candidates from the same expression compete directly. Let $\mathcal{P}\subseteq\mathcal{C}$ be the positive set and $\mathcal{N}$ the negative set used for ranking (all non-positive candidates, or only those with $J_i\le\theta_{\mathrm{neg}}$, depending on the dataset). We optimize a listwise objective over the candidates of each expression:
\begin{equation}
\mathcal{L}_{\mathrm{list}}
=
\log\!\!\sum_{j\in\mathcal{P}\cup\mathcal{N}}\!\!\exp(s_j/\kappa)
-
\log\sum_{j\in\mathcal{P}}\exp(s_j/\kappa).
\end{equation}
where $\kappa$ is a temperature. This loss increases the score mass assigned to valid target hypotheses relative to the competing candidates. If no candidate reaches $\theta_{\mathrm{pos}}$, the highest-IoU candidate is used as a fallback positive so that difficult low-recall examples still contribute ranking supervision.

\paragraph{Hard-negative ranking.}
Listwise competition is complemented by an explicit margin on confusing negatives. Since $s_i$ is bounded in $[0,1]$, this term uses the unnormalized Yes/No logit gap
\begin{equation}
 r_i=\ell_i^{\mathrm{yes}}-\ell_i^{\mathrm{no}} .
\end{equation}
Let $r^{+}$ be the positive reference, i.e., the maximum (or mean) of $r$ over $\mathcal{P}$. A negative is violating when $r_j>r^{+}-\delta$, and
\begin{equation}
\mathcal{L}_{\mathrm{hard}}
=
\frac{1}{|\mathcal{H}|}
\sum_{j\in\mathcal{H}}
\operatorname{softplus}(r_j-r^{+}+\delta),
\end{equation}
where $\mathcal{H}\subseteq\mathcal{N}$ contains the violating negatives. The loss concentrates training on candidates that remain competitive with the positive reference instead of treating all background hypotheses equally.

\paragraph{Mask-aware feature alignment.}
Candidate marking defines the anchor visually, but the verifier can still over-rely on contextual objects. We therefore regularize the hidden states of the last language-model layer. All non-visual prompt tokens, i.e., the instruction and the expression, are averaged into a text vector $t$. For the visual token $n$ of frame $f$,
\begin{equation}
 a_n=\cos(h_n,t),
 \qquad
 p_n=
 \frac{\exp(a_n/\tau_a)}{\sum_{j\in f}\exp(a_j/\tau_a)},
\end{equation}
where the softmax runs over the tokens of the same frame. The candidate mask, resized to that frame's token grid, gives a binary indicator $b_n$, and
\begin{equation}
\mathcal{L}_{\mathrm{align}}
=
\frac{1}{F}\sum_{f}\Bigl(1-\sum_{n\in f} p_n b_n\Bigr).
\end{equation}
The term is applied to every scored candidate, positive or negative: whatever the label, the verifier should gather its expression-related evidence from the marked anchor rather than from the rest of the scene, and the Yes/No decision is left to the ranking losses. It is a hidden-state regularizer rather than direct supervision of attention maps.

The verifier objective is
\begin{equation}
\mathcal{L}_{V}
=
\mathcal{L}_{\mathrm{list}}
+
\lambda_h\mathcal{L}_{\mathrm{hard}}
+
\lambda_a\mathcal{L}_{\mathrm{align}}.
\end{equation}
Listwise competition answers \emph{which candidate should win}, hard-negative ranking focuses learning on the strongest distractors, and feature alignment regularizes \emph{where} the verifier obtains its evidence.

\subsection{Inference and Target Selection}
\label{sec:inference}
An expression may refer to more than one object (e.g., ``monkeys on the left''). In addition to $\hat c$, we therefore estimate the number of referred objects $\hat N\ge 1$ with a counting prompt answered by the same frozen Qwen3-VL backbone, with the verifier adapters disabled, on the expression and $M$ uniformly sampled frames. The prompt gives priority to explicit numerals, returns one for singular references, otherwise counts the visible objects satisfying the expression, and falls back to one when uncertain. Candidates in $\mathcal{C}(q)$ are ranked by $s_i$, the top $\min(\hat N,N_c)$ tubes form $\hat{\mathcal{S}}$, and the predicted mask at frame $t$ is the union of their masks:
\begin{equation}
Y_t=\bigcup_{\tau_i\in\hat{\mathcal{S}}} M_t^{(i)}.
\end{equation}
For single-object expressions, $\hat N=1$ and this reduces to selecting the highest-scoring candidate.
